% Ad M. Brutum 2.5 --- headnote
% ad-brutum-02-05, 14 April 43 BC, Rome

\textit{Cicero to M. Junius Brutus, written from Rome on
14 April 43 BC --- Perseus dateline \textit{Scr.\ Romae
xviii K.\ Mai.\ a.\ 711 (43)}, confirmed by the closing
subscript \textit{XVIII Kalend.\ Maias} and matching the
meta date. The letter is among the most politically
weighted of the surviving correspondence with Brutus: a
manifesto and a reproach, written in the brief window
between the news of Forum Gallorum (the battle fought on
14 April outside Mutina, in which Antony was checked but
not yet broken) and the news of Mutina proper (21 April).
The Ides of April scene Cicero describes in sections 3
and 4 --- the surprise reading-out of letters from both
Brutus and Antony, the latter styled \textit{Antonius
proconsul} --- is the political crux of the letter: an
attempt by some intermediary, here Celer Pilius, to draw
the Senate into recognising Antony's continued
proconsular standing while the war was still being
fought.}

\textit{The letter's argument runs in two arcs. Sections
1 and 2 are the apologia: Cicero recalls that he and
Brutus agreed on the end --- the liberty of the
commonwealth --- but differed on means after the Ides of
March, when Brutus chose to spare Antony and Cicero (he
says) wanted the kingship abolished along with the king;
``which choice would have been better we felt then to our
great pain, we are feeling now to our great peril.'' The
intervening sentence on Octavian --- \textit{nisi Caesari
Octaviano deus quidam illam mentem dedisset} ---
gracefully credits the young Caesar with the heaven-sent
inspiration that has saved them from Antony, with no
hint of the misgivings Brutus already entertains and
will soon urge in reply. Sections 3 and 4 are the
narrative of the Ides of April scene, told in the
clipped paragraphs Cicero reserves for political theatre:
Pilius's entrance, the reading of the two letters, the
Senate's astonishment, Cicero's own hesitation, Sestius's
intervention, and Labeo's clever forensic objection.
Section 5 is the political climax: Cicero presses
Brutus to abandon his characteristic clemency, in the
sentence that has been quoted ever since --- ``the right
place and time for clemency belong, and should belong, to
other circumstances and other seasons.'' What is at
stake, Cicero insists, is not the persons of Antony or
Dolabella but the existence of the republican order
itself: \textit{nec quicquam aliud decernitur hoc bello
nisi utrum simus necne}. Section 6 closes with a single
sentence about young Marcus, on Brutus's staff. The
register is high: it is one of the few letters in which
Cicero writes a public-political doctrine to a
correspondent on the equal footing of statesman to
statesman.}
